\chapter{Аналитический раздел}

В данном разделе представлены анализ предметной области, анализ известных нейросетевых решений, классификация и сравнительный анализ рассмотренных решений, формализации изображения и задачи.

\section{Анализ предметной области}

\textit{Признак генерации} -- это опирающийся на технологии машинного обучения целенаправленно созданный атрибут медиапродукта (видео, фотографии, аудиозаписи). Изображение, содержащее признак генерации, называется \textit{поддельным}, если имеющийся на нём достоверно установленный признак направлен на формирование у целевых аудиторий фальсифицированного образа социального субъекта с целью получения выгоды или нанесения ущерба. Такие изображения представляют серьезную угрозу для доверия к цифровому контенту~\cite{spb-uni}.

В данном разделе представляется история возникновения поддельных изображений и применение методов обнаружения признаков генерации на изображении.

\subsection{История возникновения поддельных изображений}

Исследования в области автоматизированного синтеза визуального контента начались в конце 1990-х годов. В 1997 году была представлена программа <<Video Rewrite>>, которая позволяла формировать видеоряд, где артикуляционная мимика лица соответствовала синтезированной аудиодорожке~\cite{video-rewrite}. Эти разработки носили академический характер и были доступны ограниченному кругу специалистов. Эволюция технологий продолжилась в последующие два десятилетия, включая такие проекты, как использование цифрового аватара Одри Хепберн в рекламной кампании 2014 года и омоложение лиц актёров в кинематографе, например, в фильме <<Ирландец>> 2019 года~\cite{audrey-hepburn}\cite{flux}. В данный период синтез видеоконтента в значительной степени оставался процессом, контролируемым разработчиками.

Существенный прогресс в развитии технологии произошел в 2014 году с разработкой Яном Гудфеллоу программного метода генеративно-состязательных ИНС~\cite{gan}. Данный метод основан на использовании двух конкурирующих искусственных нейронных сетей: генератора, создающего синтетические данные, и дискриминатора, задачей которого является обнаружение сгенерированных образцов среди всего набора данных. 

Широкое распространение поддельные изображения получили в 2017 году, когда пользователь платформы Reddit под псевдонимом <<Deepfakes>> опубликовал серию видеороликов с заменой лиц знаменитостей. Это событие продемонстрировало потенциал технологии глубокого обучения ИНС и вызвало значительный общественный резонанс~\cite{spb-uni}.

\subsection{Применение методов обнаружения признаков генерации на изображении}

Методы обнаружения признаков генерации на изображении используются в следующих областях:
\begin{itemize}
	\item верификация видеоматериалов и аудиозаписей государственными учреждениями и новостными агентствами для противодействия манипуляции общественным мнением и провоцированию социальной напряжённости~\cite{deep-fakes}\cite{spb-uni};
	\item биометрическая аутентификация и предотвращение мошеннических операций~\cite{meso4}\cite{capsforensics}\cite{vit};
	\item анализ цифровых доказательств для установления подлинности или факта модификации видеозаписей и фотографий, представляемых в судебном делопроизводстве~\cite{faceforensics}\cite{faceforensicspp};
	\item идентификация неправомерного использования изображений граждан в синтетическом контенте без их согласия~\cite{spb-uni};
	\item контроль целостности и установление авторства цифрового контента для защиты интеллектуальной собственности от несанкционированного использования~\cite{deepfakes-document}.
\end{itemize}

\section{Анализ известных нейросетевых решений}

История развития искусственных нейронных сетей берет начало в середине 20-го века с создания упрощенных математических моделей биологического нейрона -- \textit{искусственных нейронов}. Первой такой формальной моделью стала модель, предложенная Уорреном Мак-Каллоком и Уолтером Питтсом в 1943 году~\cite{neuron}. Данная модель представляет собой устройство с несколькими бинарными входами и одним бинарным выходом. Общая схема данной модели представлена на рисунке~\ref{img:neuron}.

\begin{figure}[H]
	\centering
	\begin{tikzpicture}[>=stealth, node distance=1.6cm]
		
		% Inputs
		\node[circle, draw, minimum size=12mm] (x1) {$x_1$};
		\node[circle, draw, below=0.9cm of x1, minimum size=12mm] (x2) {$x_2$};
		\node[below=0.45cm of x2] (dots) {$\vdots$};
		\node[circle, draw, below=0.5cm of dots, minimum size=12mm] (xn) {$x_n$};
		
		% Bias input
		\node[circle, draw, right=2.3cm of xn, minimum size=12mm] (bias) {$b$};
		
		% Sum block
		\node[draw, line width=1pt, minimum width=20mm, minimum height=18mm, right=2.2cm of x2] (sum) {$\overset{n}{\underset{i=1}{\sum}}{w_ix_i}+b$};
		
		% Activation block
		\node[draw, line width=1pt, minimum width=14mm, minimum height=18mm, right=1.25cm of sum] (act) {$f_{\theta}(s)$};
		
		% Output
		\node[circle, draw, minimum size=10mm, right=1cm of act] (a) {$\alpha$};
		
		% Arrows with weights
		\draw[->] (x1) -- node[above, xshift=0.1cm, yshift=0.1cm] {$w_1$} (sum);
		\draw[->] (x2) -- node[above, xshift=-0.1cm] {$w_2$} (sum);
		\draw[->] (xn) -- node[above, xshift=-0.25cm] {$w_n$} (sum);
		\draw[->] (bias) -- node[above, xshift=-0.4cm, yshift=-0.35cm] {$1$} (sum);
		
		% Sum to activation
		\draw[->] (sum) -- node[above] {$s$} (act);
		
		% Activation to output
		\draw[->] (act) -- (a);
		
	\end{tikzpicture}
	\caption{Общая схема искусственного нейрона Мак-Каллока-Питтса~\cite{neuron}}
	\label{img:neuron}
\end{figure}

Искусственный нейрон осуществляет вычисление взвешенной суммы~$\underset{i=1}{\overset{n}{\sum}}{x_iw_i}+b$ входных сигналов, которая затем сравнивается с некоторым пороговым значением $\theta$. Если сумма превышает порог, выход нейрона $\alpha$ принимает единичное значение, в противном случае -- нулевое~\cite{mei}.

\textit{Искусственная нейронная сеть} представляет собой совокупность взаимосвязанных искусственных нейронов. Связи между нейронами определяются системой весов, которые задают структуру сети и обеспечивают преобразование входного векторного сигнала в выходной. Упорядоченная совокупность искусственных нейронов, выполняющая одновременное и идентичное преобразование входного векторного сигнала в выходной сигнал называется \textit{слоем}. ИНС может содержать неограниченное число слоёв~\cite{rumelhart}.

Одним из первых формализованных принципов изменения весов связей между нейронами является правило Хэбба, предложенное в 1949 году~\cite{hebb}. Согласно данному принципу, изменение веса связи между двумя нейронами пропорционально произведению их активностей. Если возбуждение одного нейрона регулярно сопровождается возбуждением другого, вес связи между ними увеличивается. 

Целенаправленная настройка параметров искусственной нейронной сети на основе анализа входных данных для достижения требуемого поведения модели называется \textit{обучением}~\cite{rosenblatt}.

\subsection{Свёрточные нейронные сети}

\textit{Свёрточные искусственные нейронные сети} представляют собой разновидность многослойных ИНС, разработанную для обработки данных с сеточной структурой, таких как изображения. Данная архитектура была предложена Яном Лекуном в 1988 году и вдохновлена биологическими процессами человеческой зрительной системы~\cite{convnet}.

Ключевой концепцией свёрточных сетей является использование операции \textit{свёртки}, которая позволяет выявлять локальные пространственные зависимости в обрабатываемых данных~\cite{conv}. 

Архитектура свёрточной нейронной сети включает следующие компоненты:
\begin{enumerate}
	\item \textit{Свёрточные слои} выполняют извлечение признаков путём применения набора обучаемых фильтров к входным данным. Каждый фильтр сканирует изображение, обнаруживая определённые паттерны~\cite{conv};
	\item \textit{Подвыборочные слои} уменьшают размерность карт признаков, сохраняя наиболее значимую информацию. Это позволяет снизить вычислительную сложность~\cite{pooling};
	\item \textit{Полносвязные слои} завершают архитектуру сети, преобразуя извлечённые иерархические признаки в конечный результат~\cite{convnet}\cite{rumelhart}.
\end{enumerate}

Пример схемы свёрточной искусственной нейронной сети представлен на рисунке~\ref{img:convnet}.

\begin{figure}[H]
	\centering
	\begin{tikzpicture}[
		line width=0.8pt,
		>=Stealth,
		node distance=0.3cm and 1.5cm,
		layer/.style={draw, fill=white, rounded corners=3pt, 
			minimum width=2cm, minimum height=2cm},
		smalllayer/.style={layer, minimum width=1.2cm, minimum height=1.2cm},
		fclayer/.style={layer, minimum width=0.8cm, minimum height=0.8cm},
		convcolor/.style={fill=blue!20},
		poolcolor/.style={fill=green!20},
		fccolor/.style={fill=red!20}
		]
		
		\pgfdeclarelayer{background}
		\pgfdeclarelayer{foreground}
		\pgfsetlayers{background,main,foreground}
		
		\coordinate (top-level) at (0,2.6); % фиксированная Y-координата
		\coordinate (low-level) at (0,-2.5); % фиксированная Y-координата
		
		% Изображение (вход)
		\node[layer, fill=gray!15] (input) at (0,0) {};
		\node at (input |- top-level) {\texttt{Изображение}};
		\node at ($(input.east)+(0.7,0)$) {\huge$\rightarrow$};
		
		% Первый свёрточный слой
		\convlayer{input}{conv1}{5}
		\node[align=center] at (conv12 |- top-level) {\texttt{Свёрточный} \\ \texttt{слой}};
		\node at ($(input.east)+(5.45,0)$) {\huge$\rightarrow$};
		
		% Первый подвыборочный слой
		\poollayer{conv12}{pool1}{2}
		\node[align=center] at (pool11 |- top-level) {\texttt{Подвыборочный} \\ \texttt{слой}};
		\node at ($(input.east)+(9.4,0)$) {\huge$\rightarrow$};
		
		\fclayer{pool11}{fc1}{0}
		\node[align=center] at (fc10 |- top-level) {\texttt{Полносвязный} \\ \texttt{слой}};
		
		
	\end{tikzpicture}
	\caption{Пример схемы свёрточной ИНС}
	\label{img:convnet}
\end{figure}

\subsubsection{Meso-4}

Архитектура Meso-4 была разработана в 2018 году исследовательской группой под руководством Дариуса Афшара. Она предназначена для обнаружения дипфейков на изображениях лиц людей. Основная идея архитектуры заключается в анализе \textit{мезоскопических свойств изображений}~--~промежуточного уровня между анализом цветов пикселей и их семантическим анализом~\cite{meso4}.

Данная свёрточная ИНС состоит из последовательности четырёх блоков, каждый из которых включает свёрточные и подвыборочные слои. В свёрточных слоях используется пороговая функция активации под названием Rectified Linear Unit, имеющая следующий вид~\cite{relu}:
\begin{equation}
	ReLU(z) = \max(0,z)
\end{equation}
Для улучшения обобщающей способности применяются операция \textit{пакетной нормализации}~\cite{batch} после каждого свёрточного слоя и операция \textit{отключения нейронов}~\cite{dropout} в полносвязных слоях. Схема архитектуры Meso-4 представлена на рисунке~\ref{img:meso4}.

Обучающие наборы данных и результаты подсчёта долей верно классифицированных изображений~\cite{meso4} представлены в таблицах~\ref{tbl:meso4-datasets}~и~\ref{tbl:meso4-results} соответственно.

Основным преимуществом данной архитектуры является ее относительная компактность, включающая 27\,977 параметров, что обеспечивает меньшие вычислительные затраты и возможность быстрого обучения на ограниченных ресурсах. К недостаткам можно отнести снижение точности при сжатии входных изображений~\cite{meso4}.

\begin{table}[H]
	\centering
	\caption{Результаты тестирования Meso-4}
	\label{tbl:meso4-results}
	\begin{tabular}{|c|c|}
		\hline
		\textbf{Контрольный набор данных} & \textbf{Точность, \%} \\
		\hline
		Авторский набор данных  & 89.1 \\
		\hline
		FaceForensics~\cite{faceforensics} & 94.6 \\
		\hline
	\end{tabular}
\end{table}

\begin{figure}[H]
	\centering
	\begin{tikzpicture}[node distance=6mm]
		
		\tikzset{
			layer/.style={
				rectangle, draw, minimum width=6.3cm,
				minimum height=0.5cm, align=center,
				font=\small, fill=#1
			},
			layer/.default=white,
			blockbox/.style={
				draw, rectangle, rounded corners,
				inner sep=6pt
			},
			arrow/.style={
				-latex, thick
			}
		}
		
		% ---------- Left branch ----------
		\node[layer=gray!10] (input){\texttt{Изображение 256$\times$256$\times$3}};
		
		\node[layer=blue!20, below=0.8cm of input] (c1) {\texttt{Свёрточный 8$\times$(3$\times$3) + ReLU}};
		\node[layer=yellow!10, below=0.1cm of c1] (bn1) {\texttt{Пакетная нормализация}};
		\node[layer=green!20, below=0.1cm of bn1] (mp1) {\texttt{Подвыборочный 2$\times$2}};
		\node[blockbox, fit=(c1) (bn1) (mp1), below=of input] (block1) {};
		
		\node[layer=blue!20, below=1.05cm of mp1] (c2){\texttt{Свёрточный 8$\times$(5$\times$5) + ReLU}};
		\node[layer=yellow!10, below=0.1cm of c2] (bn2){\texttt{Пакетная нормализация}};
		\node[layer=green!20, below=0.1cm of bn2] (mp2){\texttt{Подвыборочный 2$\times$2}};
		\node[blockbox, fit=(c2) (bn2) (mp2), below=of block1] (block2) {};
		
		\node[layer=blue!20, below=1.05cm of mp2] (c3){\texttt{Свёрточный 16$\times$(5$\times$5) + ReLU}};
		\node[layer=yellow!10, below=0.1cm of c3] (bn3){\texttt{Пакетная нормализация}};
		\node[layer=green!20, below=0.1cm of bn3] (mp3){\texttt{Подвыборочный 2$\times$2}};
		\node[blockbox, fit=(c3) (bn3) (mp3), below=of block2] (block3) {};
		
		\draw[arrow] (input) -- (block1);
		\draw[arrow] (block1) -- (block2)node[midway, left, font=\small\ttfamily]{128$\times$128$\times$8};
		\draw[arrow] (block2) -- (block3)node[midway, left, font=\small\ttfamily]{64$\times$64$\times$8};
		
		% ---------- Right branch (Meso-4) ----------
		\node[layer=blue!20, right=1.05cm of c1] (c4){\texttt{Свёрточный 16$\times$(5$\times$5) + ReLU}};
		\node[layer=yellow!10, below=0.1cm of c4] (bn4){\texttt{Пакетная нормализация}};
		\node[layer=green!20, below=0.1cm of bn4] (mp4){\texttt{Подвыборочный 2$\times$2}};
		\node[blockbox, fit=(c4) (bn4) (mp4), right=of block1] (block4) {};
		
		\node[layer=orange!20, below=1.05cm of mp4] (d5){\texttt{Отключение нейронов 0.5}};
		\node[layer=red!20, below=0.1cm of d5] (f5){\texttt{Полносвязный 16}};
		\node[blockbox, fit=(d5) (f5), below=of block4] (block5) {};
		
		\node[layer=orange!20, below=1.05cm of f5] (d6){\texttt{Отключение нейронов 0.5}};
		\node[layer=red!20, below=0.1cm of d6] (f6){\texttt{Полносвязный 1 + Sigmoid}};
		\node[blockbox, fit=(d6) (f6), below=of block5] (block6) {};
		
		\node[layer=gray!10, below=0.8cm of f6] (output){\texttt{Результат}};
		
		
		\coordinate (p1) at ($(block3.south)+(0,-0.3)$);   % вниз
		\coordinate (p2) at ($(p1)+(3.7,0)$);              % вправо
		\coordinate (p3) at ($(p2)+(0,10.25)$); 			   % вверх 
		\coordinate (p4) at ($(p3)+(3.65,0)$); 			   % вправо 
		\draw[arrow] 
		(block3.south) -- 
		(p1) -- 
		(p2) -- 
		(p3) -- 
		(p4) --
		(block4.north)
		node[midway, right, font=\small\ttfamily]{32$\times$32$\times$16};
		\draw[arrow] (block4) -- (block5)node[midway, right, font=\small\ttfamily]{1024};
		\draw[arrow] (block5) -- (block6)node[midway, right, font=\small\ttfamily]{16};
		\draw[arrow] (block6) -- (output)node[midway, right, font=\small\ttfamily]{1};	
	\end{tikzpicture}
	\caption{Схема архитектуры Meso-4}
	\label{img:meso4}
\end{figure}


\subsubsection{MesoInception-4}

Архитектура MesoInception-4, предложенная в 2018 году Дариусом Афшаром, представляет собой модификацию архитектуры Meso-4~\cite{meso4}. Её ключевая особенность заключается в наличии \textit{слоёв многомасштабной свёртки}~\cite{inception} на начальных этапах сети.

Идея слоя многомасштабной свёртки заключается в применении параллельных операций свёртки с ядрами различного размера. Данное решение позволяет анализировать признаки изображения на разных уровнях детализации в рамках одного вычислительного блока. Схемы архитектуры данного слоя и архитектуры MesoInception-4 представлены на рисунках~\ref{img:mesoinception4-inception}~и~\ref{img:mesoinception4} соответственно.

\begin{table}[h]
	\centering
	\caption{Параметры слоёв многомасштабной свёртки архитектуры MesoInception-4}
	\label{tbl:mesoinception4-inception-params}
	\begin{tabular}{|c|c|c|c|c|}
		\hline
		\textbf{Номер слоя} & \textbf{A} & \textbf{B} & \textbf{C} & \textbf{D} \\
		\hline
		1 & 1 & 4 & 4 & 1 \\
		\hline
		2 & 1 & 4 & 4 & 2 \\
		\hline
	\end{tabular}
\end{table}

Параметры слоёв многомасштабной свёртки, обучающие наборы данных и результаты подсчёта долей верно классифицированных изображений~\cite{meso4} представлены в таблицах~\ref{tbl:mesoinception4-inception-params},~\ref{tbl:meso4-datasets}~и~\ref{tbl:mesoinception4-results} соответственно.

\begin{figure}[h]
	\centering
	\begin{tikzpicture}[node distance=6mm]
		
		\tikzset{
			layer/.style={
				rectangle, draw, minimum width=6.6cm,
				minimum height=0.5cm, align=center,
				font=\small, fill=#1
			},
			layer/.default=white,
			blockbox/.style={
				draw, rectangle, rounded corners,
				inner sep=6pt
			},
			arrow/.style={
				-latex, thick
			}
		}
		
		% ---------- Left branch ----------
		\node[layer=gray!10] (input){\texttt{Входные данные L$\times$L$\times$P}};
		
		\node[layer=blue!10, below=0.8cm of input] (conv1x1-1) {\texttt{Свёрточный A$\times$(1$\times$1)}};
		\node[blockbox, fit=(conv1x1-1)] (branch1) {};
		
		\node[layer=blue!10, below=1.05cm of conv1x1-1] (conv1x1-2) {\texttt{Свёрточный B$\times$(1$\times$1)}};
		\node[layer=blue!20, below=0.1cm of conv1x1-2] (dilated3x3-1) {\texttt{Свёрточный B$\times$(3$\times$3)}};
		\node[blockbox, fit=(conv1x1-2) (dilated3x3-1), below=of branch1] (branch2) {};
		
		\node[layer=blue!10, below=1.05cm of dilated3x3-1] (conv1x1-3) {\texttt{Свёрточный C$\times$(1$\times$1)}};
		\node[layer=blue!20, below=0.1cm of conv1x1-3] (dilated3x3-2) {\texttt{Свёрточный C$\times$(3$\times$3)}};
		\node[blockbox, fit=(conv1x1-3) (dilated3x3-2), below=of branch2] (branch3) {};
		
		\node[layer=blue!10, below=1.05cm of dilated3x3-2] (conv1x1-4) {\texttt{Свёрточный D$\times$(1$\times$1)}};
		\node[layer=blue!20, below=0.1cm of conv1x1-4] (dilated3x3-3) {\texttt{Свёрточный D$\times$(3$\times$3)}};
		\node[blockbox, fit=(conv1x1-4) (dilated3x3-3), below=of branch3] (branch4) {};
		
		\coordinate (p1) at ($(input.south)+(0,-0.3)$);
		\coordinate (p2) at ($(p1)+(-4,0)$);
		
		\draw[arrow]
		(input.south) --
		(p1) --
		(p2) --
		($(branch1.west)-(0.475,0)$) --
		(branch1.west);
		
		\draw[arrow]
		(input.south) --
		(p1) --
		(p2) --
		($(branch2.west)-(0.475,0)$) --
		(branch2.west);
		
		\draw[arrow]
		(input.south) --
		(p1) --
		(p2) --
		($(branch3.west)-(0.475,0)$) --
		(branch3.west);
		
		\draw[arrow]
		(input.south) --
		(p1) --
		(p2) --
		($(branch4.west)-(0.475,0)$) --
		(branch4.west);
		
		% ---------- Right branch ----------
		\node[layer=gray!10, right=1.25cm of input] (output){\texttt{Выходные данные L$\times$L$\times$(A+B+C+D)}};
		
		\node[layer=pink!50, right=1.3625cm of conv1x1-1] (conc) {\texttt{Конкатенация}};
		\node[blockbox, fit=(conc)] (branch5) {};
		
		% Создаём промежуточные координаты для выравнивания стрелок
		\coordinate (vertical-line) at ($(branch5.west)-(0.5,0)$);
		
		% Выравниваем стрелки по вертикали
		\draw[arrow] (branch1.east) -- (vertical-line |- branch5.west) -- (branch5.west);
		\draw[arrow] (branch2.east) -- (branch2.east -| vertical-line) -- (vertical-line |- branch5.west) -- (branch5.west);
		\draw[arrow] (branch3.east) -- (branch3.east -| vertical-line) -- (vertical-line |- branch5.west) -- (branch5.west);
		\draw[arrow] (branch4.east) -- (branch4.east -| vertical-line) -- (vertical-line |- branch5.west) -- (branch5.west);
		\draw[arrow] (branch5.north) -- (output.south);
		
	\end{tikzpicture}
	\caption{Схема архитектуры слоя многомасштабной свёртки, параметризуемой с помощью $A, B, C, D \in \mathbb{N}$}
	\label{img:mesoinception4-inception}
\end{figure}

Главным преимуществом архитектуры MesoInception-4 является способность комбинировать признаки различного масштаба благодаря использованию слоёв многомасштабной свёртки, что привело к небольшому повышению точности по сравнению с Meso-4, при практически неизменной вычислительной сложности. Недостатком является чувствительность к сжатию входных изображений.

\begin{table}[H]
	\centering
	\caption{Результаты тестирования MesoInception-4}
	\label{tbl:mesoinception4-results}
	\begin{tabular}{|c|c|}
		\hline
		\textbf{Контрольный набор данных} & \textbf{Точность, \%} \\
		\hline
		Авторский набор данных & 91.7 \\
		\hline
		FaceForensics & 96.8 \\
		\hline
	\end{tabular}
\end{table}

\begin{figure}[H]
	\centering
	\begin{tikzpicture}[node distance=6mm]
		
		\tikzset{
			layer/.style={
				rectangle, draw, minimum width=6.3cm,
				minimum height=0.5cm, align=center,
				font=\small, fill=#1
			},
			layer/.default=white,
			blockbox/.style={
				draw, rectangle, rounded corners,
				inner sep=6pt
			},
			arrow/.style={
				-latex, thick
			}
		}
		
		% ---------- Left branch ----------
		\node[layer=gray!10] (input){\texttt{Изображение 256$\times$256$\times$3}};
		
		\node[layer=blue!20, below=0.8cm of input] (mp1) {\texttt{Многомасштабная свёртка}};
		\node[blockbox, fit=(mp1), below=of input] (block1) {};
		
		\node[layer=blue!20, below=1.05cm of mp1] (mp2){\texttt{Многомасштабная свёртка}};
		\node[blockbox, fit=(mp2), below=of block1] (block2) {};
		
		\node[layer=blue!20, below=1.05cm of mp2] (c3){\texttt{Свёрточный 16$\times$(5$\times$5) + ReLU}};
		\node[layer=yellow!10, below=0.1cm of c3] (bn3){\texttt{Пакетная нормализация}};
		\node[layer=green!20, below=0.1cm of bn3] (mp3){\texttt{Подвыборочный 2$\times$2}};
		\node[blockbox, fit=(c3) (bn3) (mp3), below=of block2] (block3) {};
		
		\draw[arrow] (input) -- (block1);
		\draw[arrow] (block1) -- (block2)node[midway, left, font=\small\ttfamily]{};
		\draw[arrow] (block2) -- (block3)node[midway, left, font=\small\ttfamily]{};
		
		% ---------- Right branch ----------
		\node[layer=blue!20, right=1.05cm of mp1, yshift=0.8cm] (c4){\texttt{Свёрточный 16$\times$(5$\times$5) + ReLU}};
		\node[layer=yellow!10, below=0.1cm of c4] (bn4){\texttt{Пакетная нормализация}};
		\node[layer=green!20, below=0.1cm of bn4] (mp4){\texttt{Подвыборочный 2$\times$2}};
		\node[blockbox, fit=(c4) (bn4) (mp4), right=of block1] (block4) {};
		
		\node[layer=orange!20, below=1.05cm of mp4] (d5){\texttt{Отключение нейронов 0.5}};
		\node[layer=red!20, below=0.1cm of d5] (f5){\texttt{Полносвязный 16}};
		\node[blockbox, fit=(d5) (f5), below=of block4] (block5) {};
		
		\node[layer=orange!20, below=1.05cm of f5] (d6){\texttt{Отключение нейронов 0.5}};
		\node[layer=red!20, below=0.1cm of d6] (f6){\texttt{Полносвязный 1 + Sigmoid}};
		\node[blockbox, fit=(d6) (f6), below=of block5] (block6) {};
		
		\node[layer=gray!10, below=0.8cm of f6] (output){\texttt{Результат}};
		
		
		\coordinate (p1) at ($(block3.south)+(0,-0.3)$);   % вниз
		\coordinate (p2) at ($(p1)+(3.7,0)$);              % вправо
		\coordinate (p3) at ($(p2)+(0,7.75)$); 			   % вверх 
		\coordinate (p4) at ($(p3)+(3.65,0)$); 			   % вправо 
		\draw[arrow] 
		(block3.south) -- 
		(p1) -- 
		(p2) -- 
		(p3) -- 
		(p4) --
		(block4.north)
		node[midway, right, font=\small\ttfamily]{};
		\draw[arrow] (block4) -- (block5)node[midway, right, font=\small\ttfamily]{1024};
		\draw[arrow] (block5) -- (block6)node[midway, right, font=\small\ttfamily]{16};
		\draw[arrow] (block6) -- (output)node[midway, right, font=\small\ttfamily]{1};	
	\end{tikzpicture}
	\caption{Схема архитектуры MesoInception-4}
	\label{img:mesoinception4}
\end{figure}

\subsection{Капсульные нейронные сети}

\textit{Капсульные искусственные нейронные сети} представляют собой архитектуру, предложенную в 2017 году Джеффри Хинтоном~\cite{capsnn}, целью которой является преодоление ограничений свёрточных ИНС в области инвариантности к пространственным преобразованиям и анализа иерархических отношений между частями объектов. Пример схемы капсульной ИНС представлен на рисунке~\ref{img:capsnet}.

\begin{figure}[H]
	\centering
	\begin{tikzpicture}[
		line width=0.8pt,
		>=Stealth,
		node distance=0.3cm and 1.5cm,
		layer/.style={draw, fill=white, rounded corners=3pt, 
			minimum width=2cm, minimum height=2cm},
		smalllayer/.style={layer, minimum width=1.2cm, minimum height=1.2cm},
		fclayer/.style={layer, minimum width=0.8cm, minimum height=0.8cm},
		convcolor/.style={fill=blue!20},
		poolcolor/.style={fill=green!20},
		fccolor/.style={fill=red!20},
		capsulecolor/.style={fill=orange!20}
		]
		
		\pgfdeclarelayer{background}
		\pgfdeclarelayer{foreground}
		\pgfsetlayers{background,main,foreground}
		
		\coordinate (top-level) at (0,2.75); % фиксированная Y-координата
		\coordinate (low-level) at (0,-2.5); % фиксированная Y-координата
		
		% Изображение (вход)
		\node[layer, fill=gray!15] (input) at (0,0) {};
		\node at (input |- top-level) {\texttt{Изображение}};
		\node at ($(input.east)+(0.7,0)$) {\huge$\rightarrow$};
		
		% Свёрточный слой
		\convlayer{input}{conv1}{5}
		\node[align=center] at (conv12 |- top-level) {\texttt{Свёрточный} \\ \texttt{слой}};
		\node at ($(input.east)+(5.3,0)$) {\huge$\rightarrow$};
		
		% капсулы
		\capsulelayer{conv12}{capsule1}{5}
		\node[align=center] at (capsule12 |- top-level) {\texttt{Капсульный} \\ \texttt{слой}};
		\node at ($(input.east)+(10.15,0)$) {\huge$\rightarrow$};
		
		% Выходные капсулы (классы)
		\fclayer{capsule12}{fc1}{4}
		\node[align=center] at (fc11 |- top-level) {\texttt{Полносвязный} \\ \texttt{слой}};
		
	\end{tikzpicture}
	\caption{Пример схемы капсульной ИНС}
	\label{img:capsnet}
\end{figure}

Основной единицей данной архитектуры является \textit{капсула} -- группа нейронов, активность которых отражает вероятность обнаружения определённого объекта или его части, а также его параметры, такие как положение, ориентация, масштаб и иные пространственные свойства. В отличие от скалярных выходов нейронов в свёрточных сетях или перцептронах, выход капсулы является вектором, длина которого определяет вероятность обнаружения объекта, а компоненты содержат значения его параметров. Пример схемы структуры капсулы представлена на рисунке~\ref{img:capsule-structure}.

\begin{figure}[h]
	\centering
	\begin{tikzpicture}[
		line width=0.8pt,
		>=Stealth,
		neuron/.style={draw, circle, minimum size=1cm},
		capsule/.style={draw, fill=orange!20, rounded corners=5pt, 
			minimum width=2cm, minimum height=8cm}
		]
		
		\coordinate (top-level) at (0,5);
		
		% Капсула
		\node[capsule] (capsule) at (0,0) {};
		
		% Нейроны внутри капсулы с подписями
		\node[neuron, fill=yellow!15] (n1) at (0,3) {$s_{j1}$};
		\node[neuron, fill=yellow!15] (n2) at (0,1.5) {$s_{j2}$};
		\node[neuron, fill=yellow!15] (n3) at (0,0.0) {$s_{j3}$};
		\node[neuron, fill=yellow!15] (n4) at (0,-1.5) {$s_{j4}$};
		\node[neuron, fill=yellow!15] (n5) at (0,-3) {$s_{j5}$};
		
		% Входы
		\node[neuron, fill=blue!15] [left=3.5cm of n1] (in1) {$\vec{v_{i1}}$};
		\node[neuron, fill=blue!15] [left=3.5cm of n2] (in2) {$\vec{v_{i2}}$};
		\node[neuron, fill=blue!15] [left=3.5cm of n3] (in3) {$\vec{v_{i3}}$};
		\node[neuron, fill=blue!15] [left=3.5cm of n4] (in4) {$\vec{v_{i4}}$};
		\node[neuron, fill=blue!15] [left=3.5cm of n5] (in5) {$\vec{v_{i5}}$};
		
		% Выходной вектор
		\node[neuron, fill=red!15] [right=3.5cm of capsule] (output) {$\vec{v}_j$};
		
		\coordinate (nn1) at ($(n1.west)-(0.4,0)$);
		\coordinate (nn2) at ($(n2.west)-(0.4,0)$);
		\coordinate (nn3) at ($(n3.west)-(0.4,0)$);
		\coordinate (nn4) at ($(n4.west)-(0.4,0)$);
		\coordinate (nn5) at ($(n5.west)-(0.4,0)$);
		
		\foreach \i in {1,...,5}
		\foreach \j in {1,...,5}
		\draw[->] (in\i) -- (nn\j) -- (n\j);
		
		
		\draw[->] (n1) -- (output);
		\draw[->] (n2) -- (output);
		\draw[->] (n3) -- (output);
		\draw[->] (n4) -- (output);
		\draw[->] (n5) -- (output);
		
		% Поясняющие надписи
		\node[align=center] at (in1 |- top-level) {\texttt{Входные векторы}};
		\node[align=center] at (capsule |- top-level) {\texttt{Группа нейронов}};
		\node[align=center] at (output |- top-level) {\texttt{Выходной вектор}};
		
	\end{tikzpicture}
	\caption{Пример схемы структуры капсулы}
	\label{img:capsule-structure}
\end{figure}

Алгоритмом, обеспечивающим взаимодействие между капсулами, является алгоритм \textit{динамической маршрутизации}~\cite{capsnn}, определяющий степень связи между капсулами смежных слоёв. Формально, выходной вектор капсулы вычисляется с помощью нелинейной \textit{функции сжатия}, которая сохраняет ориентацию вектора и нормализует его длину:
\begin{equation}
	\vec{v}_j = \frac{\|\vec{s}_j\|^2}{1 + \|\vec{s}_j\|^2} \cdot \frac{\vec{s}_j}{\|\vec{s}_j\|},
\end{equation}
где $\vec{s_j}$ -- вектор выходных значений нейронов капсулы.

Архитектура капсульных сетей отличается от свёрточных сетей наличием капсульных слоёв, меньшим количеством или полным отсутствием подвыборочных слоёв и большей подверженности воздействию шумов на входных данных~\cite{capsnn}. 

\subsubsection{Capsule-Forensics}

Капсульная сеть Capsule-Forensics была представлена в 2018 году авторами Х. Нгуеном, Джуничи Ямагиши и Исао Эчизеном~\cite{capsforensics}. Это одна из первых архитектур, применяющих капсульные сети для обнаружения дипфейков на изображениях.

Ключевая идея данной архитектуры заключается в использовании предобученных свёрточных слоёв сети VGG-19~\cite{vgg19} для извлечения признаков входных изображений и их последующей обработки капсульными слоями. Схема архитектуры Capsule-Forensics представлена на рисунке~\ref{img:capsule-forensics}.

\begin{figure}[H]
	\centering
	\begin{tikzpicture}[node distance=6mm]
		
		\tikzset{
			layer/.style={
				rectangle, draw, minimum width=6.5cm,
				minimum height=0.5cm, align=center,
				font=\small, fill=#1
			},
			layer/.default=white,
			blockbox/.style={
				draw, rectangle, rounded corners,
				inner sep=6pt
			},
			arrow/.style={
				-latex, thick
			}
		}
		
		% Вход
		\node[layer=gray!10] (input){\texttt{Изображение 224$\times$224$\times$3}};
		
		% VGG-19 Feature Extractor
		\node[layer=blue!20, below=0.8cm of input] (vgg-conv1) {\texttt{Свёрточный 64$\times$(3$\times$3) + ReLU}};
		\node[layer=blue!20, below=0.1cm of vgg-conv1] (vgg-conv2) {\texttt{Свёрточный 64$\times$(3$\times$3) + ReLU}};
		\node[layer=green!20, below=0.1cm of vgg-conv2] (vgg-pool1) {\texttt{Подвыборочный 2$\times$2}};
		\node[blockbox, fit=(vgg-conv1) (vgg-conv2) (vgg-pool1), below=of input] (block1) {};
		
		\node[layer=blue!20, below=1.05cm of vgg-pool1] (vgg-conv3) {\texttt{Свёрточный 128$\times$(3$\times$3) + ReLU}};
		\node[layer=blue!20, below=0.1cm of vgg-conv3] (vgg-conv4) {\texttt{Свёрточный 128$\times$(3$\times$3) + ReLU}};
		\node[layer=green!20, below=0.1cm of vgg-conv4] (vgg-pool2) {\texttt{Подвыборочный 2$\times$2}};
		\node[blockbox, fit=(vgg-conv3) (vgg-conv4) (vgg-pool2), below=of block1] (block2) {};
		
		\node[layer=blue!20, below=1.05cm of vgg-pool2] (vgg-conv5) {\texttt{Свёрточный 256$\times$(3$\times$3) + ReLU}};
		\node[layer=blue!20, below=0.1cm of vgg-conv5] (vgg-conv6) {\texttt{Свёрточный 256$\times$(3$\times$3) + ReLU}};
		\node[layer=blue!20, below=0.1cm of vgg-conv6] (vgg-conv7) {\texttt{Свёрточный 256$\times$(3$\times$3) + ReLU}};
		\node[layer=blue!20, below=0.1cm of vgg-conv7] (vgg-conv8) {\texttt{Свёрточный 256$\times$(3$\times$3) + ReLU}};
		\node[layer=green!20, below=0.1cm of vgg-conv8] (vgg-pool3) {\texttt{Подвыборочный 2$\times$2}};
		\node[blockbox, fit=(vgg-conv5) (vgg-conv6) (vgg-conv7) (vgg-conv8) (vgg-pool3), below=of block2] (block3) {};
		
		% Primary Capsules
		\node[layer=orange!20, right=1.05cm of vgg-conv1, yshift=-0.8cm] (primary1) {\texttt{Капсульный 3}};
		\node[blockbox, fit=(primary1), right=of block1] (caps-block) {};
		
		% Output Capsules
		\node[layer=orange!20, below=1.05cm of primary1] (output1) {\texttt{Капсульный 2}};
		\node[blockbox, fit=(output1), below=of caps-block] (output-block) {};
		
		\node[layer=red!20, below=1.05cm of output1] (fc1) {\texttt{Полносвязный 1 + Softmax}};
		\node[blockbox, fit=(output1), below=of output-block] (fc-block) {};
		
		% Final Output
		\node[layer=gray!10, below=0.8cm of fc1] (final){\texttt{Результат}};
		
		% Соединения
		\draw[arrow] (input) -- (block1);
		\draw[arrow] (block1) -- (block2)node[midway, left, font=\small\ttfamily]{224$\times$224$\times$64};
		\draw[arrow] (block2) -- (block3)node[midway, left, font=\small\ttfamily]{112$\times$112$\times$128};
		
		\coordinate (p1) at ($(block3.south)+(0,-0.35)$);   % вниз
		\coordinate (p2) at ($(p1)+(3.75,0)$);              % вправо
		\coordinate (p3) at ($(p2)+(0,11.15)$); 			   % вверх 
		\coordinate (p4) at ($(p3)+(3.85,0)$);
		\draw[arrow] 
		(block3) -- 
		(p1) --
		(p2) --
		(p3) --
		(p4) --
		(caps-block)node[midway, right, font=\small\ttfamily]{56$\times$56$\times$256};
		
		\draw[arrow] (caps-block) -- (output-block);
		\draw[arrow] (output-block) -- (fc-block);
		\draw[arrow] (fc-block) -- (final);
		
	\end{tikzpicture}
	\caption{Схема архитектуры Capsule-Forensics}
	\label{img:capsule-forensics}
\end{figure}

Наличие капсульных слоёв позволяет данной сети лучше определять пространственные иерархии и отношения между частями изображения по сравнению со свёрточными сетями~\cite{capsnn}. Недостатками можно считать относительно сложную архитектуру, вычислительную сложность алгоритма динамической маршрутизации, а также зависимость от предобученной модели VGG-19~\cite{capsforensics}.

В выходном полносвязном слое используется пороговая функция активации под названием Softmax, которая имеет следующий вид~\cite{softmax}:
\begin{equation}
	Softmax(\vec{z})_i = \frac{e^{z_i}}{\sum_{j=1}^{K} e^{z_j}}, \quad i=\overline{1,K},
\end{equation}
где
\begin{itemize}[label=---]
	\item $\vec{z} = (z_1, z_2, \dots, z_K)$ -- входной вектор;
	\item $K$ -- количество параметров.
\end{itemize}

Обучающий набор данных рассматриваемой ИНС представлен в виде FaceForensics. Результаты тестирования представлены в таблице~\ref{tbl:capsule-forensics-results}~\cite{capsforensics}.

\begin{table}[H]
	\centering
	\caption{Результаты тестирования Capsule-Forensics}
	\label{tbl:capsule-forensics-results}
	\begin{tabular}{|c|c|}
		\hline
		\textbf{Контрольный набор данных} & \textbf{Точность, \%} \\
		\hline
		FaceForensics & 99.37 \\
		\hline
	\end{tabular}
\end{table}

\subsection{Трансформеры}

\textit{Трансформеры} представляют собой архитектуру искусственных нейронных сетей, предназначенную для обработки последовательностей данных. Данная архитектура была предложена в 2017 году коллективом авторов из Google Brain~\cite{attention} и была применена в задаче обработки языка.

Особенностью трансформеров является использование \textit{механизма внимания}~\cite{attention}, который позволяет оценивать взаимное влияние всех элементов последовательности данных друг на друга и представлять их в форме, учитывающей это влияние.

\begin{figure}[H]
	\centering
	\begin{tikzpicture}[
		line width=0.8pt,
		>=Stealth,
		node distance=0.3cm and 1.5cm,
		layer/.style={draw, fill=white, rounded corners=3pt, 
			minimum width=2cm, minimum height=2cm},
		smalllayer/.style={layer, minimum width=1.2cm, minimum height=1.2cm},
		fclayer/.style={layer, minimum width=0.8cm, minimum height=0.8cm},
		convcolor/.style={fill=blue!20},
		poolcolor/.style={fill=green!20},
		fccolor/.style={fill=red!20},
		attentioncolor/.style={fill=yellow!20}
		]
		
		\pgfdeclarelayer{background}
		\pgfdeclarelayer{foreground}
		\pgfsetlayers{background,main,foreground}
		
		\coordinate (top-level) at (0,2.75); % фиксированная Y-координата
		
		% Изображение (вход)
		\node[layer, fill=gray!15] (input) at (0,0) {};
		\node[align=center] at (input |- top-level) {\texttt{Входные} \\ \texttt{данные}};
		\node at ($(input.east)+(0.9,0)$) {\huge$\rightarrow$};
		
		\attentionlayer{input}{attention1}{6}
		\node[align=center] at (attention12 |- top-level) {\texttt{Слой} \\ \texttt{внимания}};
		\node at ($(input.east)+(6.1,0)$) {\huge$\rightarrow$};
		
		\fclayer{attention13}{fc1}{4}
		\node[align=center] at (fc11 |- top-level) {\texttt{Полносвязный} \\ \texttt{слой}};
		
	\end{tikzpicture}
	\caption{Пример схемы трансформера}
	\label{img:transformer}
\end{figure}

Пусть задана последовательность векторов $\vec{x}_1, \vec{x}_2, \ldots, \vec{x}_n \in \mathbb{R}^d$, где $n~\in~\mathbb{N}$~--~длина последовательности, $d \in \mathbb{N}$~--~размерность векторного представления элемента. Для каждого элемента последовательности вычисляется следующий набор векторов~\cite{attention}:
\begin{enumerate}
	\item \textit{Запрос} $\vec{q}_i$ -- вектор, представляющий текущий элемент в роли <<искателя>> контекста от других элементов
	\begin{equation}
		\vec{q}_i = W^Q \vec{x}_i,
	\end{equation}
	где $W^Q$ -- обучаемая матрица параметров запросов;
	
	\item \textit{Ключ} $\vec{k}_i$ -- вектор, представляющий элемент в роли <<источника>> контекста для других элементов
	\begin{equation}
		\vec{k}_i = W^K \vec{x}_i,
	\end{equation}
	где $W^K$ -- обучаемая матрица параметров ключей;
	
	\item \textit{Значение} $\vec{v}_i$ -- вектор, содержащий фактическую информацию об элементе для агрегирования
	\begin{equation}
		\vec{v}_i = W^V \vec{x}_i,
	\end{equation}
	где $W^V$ -- обучаемая матрица параметров значений.
\end{enumerate}

Затем для каждого элемента $\vec{x}_i$ вычисляется степень его влияния на все остальные элементы последовательности. Для этого определяется скалярное произведение вектора запроса $\vec{q}_i$ на векторы ключей всех элементов $\vec{k}_j$ с последующим применением функции Softmax:
\begin{equation}
	\alpha_{ij} = \frac{\exp\left( \frac{\vec{q}_i \cdot \vec{k}_j}{\sqrt{d}} \right)}{\sum_{m=1}^n \exp\left( \frac{\vec{q}_i \cdot \vec{k}_m}{\sqrt{d}} \right)},
\end{equation}
где $\alpha_{ij} \in [0,1]$ -- степень влияния элемента $\vec{x}_j$ на элемент $\vec{x}_i$.

Выходной вектор $\vec{z}_i$ для элемента $\vec{x}_i$, учитывающий влияние других элементов,  представляет собой взвешенную сумму векторов значений~\cite{attention}:
\begin{equation}
	\vec{z}_i = \sum_{j=1}^n \alpha_{ij} \vec{v}_j
\end{equation}

В результате механизма внимания формируется новая последовательность контекстуализированных векторов $\vec{z}_1, \vec{z}_2, \ldots, \vec{z}_n$. Каждый вектор $\vec{z}_i$ содержит информацию не только об исходном элементе $\vec{x}_i$, но и о всех элементах последовательности, с которыми он значимо взаимодействует~\cite{attention}.

Трансформер может быть адаптирован для работы с изображениями. Для этого изображение разбивается на непересекающиеся фрагменты фиксированного размера. Каждый фрагмент линейно преобразуется в векторное представление, после чего получается последовательность векторов $\vec{x}_1, \vec{x}_2, \ldots, \vec{x}_n$, где $n$ -- количество фрагментов. К этой последовательности применяется стандартный механизм работы трансформера~\cite{vit}.

К преимуществам обработки изображений трансформерами относятся восприятие глобального контекста за счёт механизма внимания и  возможность параллельной обработки всех элементов последовательности. Недостатками являются квадратичная вычислительная сложность по отношению к количеству фрагментов и относительно большое количество параметров~\cite{attention}.

\subsubsection{Vision Transformer}

Архитектура Vision Transformer была представлена в 2020 году исследовательской командой Google Brain. Её ключевой особенностью является адаптация архитектуры трансформера к задаче компьютерного зрения~\cite{vit}.

Основная идея данной архитектуры заключается в разделении изображения на квадратные фрагменты $16\times16$, последующем создании векторных представлений этих фрагментов и применении стандартного механизма работы для получившейся последовательности векторов. Векторным представлением фрагмента является набор из $16\times16\times3$ элементов, компоненты которого соответствуют интенсивностям пикселей фрагмента. Схема Vision Transformer представлена на рисунке~\ref{img:vit}.

Преимуществом данной архитектуры является высокая обобщающая способность. Недостатком же является относительно большая вычислительная сложность и количество параметров -- 84 миллиона~\cite{vit}.

\begin{figure}[H]
	\centering
	\begin{tikzpicture}[node distance=6mm]
		
		\tikzset{
			layer/.style={
				rectangle, draw, minimum width=6.3cm,
				minimum height=0.5cm, align=center,
				font=\small, fill=#1
			},
			layer/.default=white,
			blockbox/.style={
				draw, rectangle, rounded corners,
				inner sep=6pt
			},
			blockbox-dashed/.style={
				draw, rectangle, rounded corners,
				inner sep=6pt,
				dashed
			},
			arrow/.style={
				-latex, thick
			}
		}
		
		% ---------- Left branch ----------
		\node[layer=gray!10] (input){\texttt{Изображение 256$\times$256$\times$3}};
		
		\node[layer=yellow!10, below=0.8cm of input] (bn1) {\texttt{Пакетная нормализация}};
		\node[layer=yellow!20, below=0.1cm of bn1] (c1) {\texttt{Слой внимания 16}};
		\node[layer=yellow!10, below=0.1cm of c1] (bn2) {\texttt{Пакетная нормализация}};
		\node[layer=red!20, below=0.1cm of bn2] (mp1) {\texttt{Полносвязный 5120 + GELU}};
		\node[blockbox-dashed, fit=(bn1) (c1) (bn2) (mp1), below=of input] (block1) {};
		\node[align=center, right=0.2cm of block1] {\texttt{$\times$ 16}};
		
		\node[layer=red!20, below=1.05cm of mp1] (mp2) {\texttt{Полносвязный 2 + Softmax}};
		\node[blockbox, fit=(mp2), below=of block1] (block2) {};
		
		\node[layer=gray!10, below=0.8cm of mp2] (output){\texttt{Результат}};
		
		\draw[arrow] (input) -- (block1);
		\draw[arrow] (block1) -- (block2);
		\draw[arrow] (block2) -- (output);
		
	\end{tikzpicture}
	\caption{Схема архитектуры Vision Transformer}
	\label{img:vit}
\end{figure}

В большинстве полносвязных слоёв архитектуры используется функция активации под названием Gaussian Error Linear Unit, имеющая следующий вид~\cite{gelu}:
\begin{equation}
	GELU(x) = 0.5x \left( 1 + \tanh\left[ \sqrt{\frac{2}{\pi}} \left( x + 0.044715x^3 \right) \right] \right)
\end{equation}

Обучающий набор данных рассматриваемой ИНС представлен в виде FaceForensics. Результаты тестирования представлены в таблице~\ref{tbl:vit-results}~\cite{dd-vit}.

\begin{table}[H]
	\centering
	\caption{Результаты тестирования Vision Transformer}
	\label{tbl:vit-results}
	\begin{tabular}{|c|c|}
		\hline
		\textbf{Контрольный набор данных} & \textbf{Точность, \%} \\
		\hline
		FaceForensics & 97.00 \\
		\hline
	\end{tabular}
\end{table}

\section{Классификация рассмотренных решений}

В данном разделе представляются критерии и проведение классификации рассмотренных решений.

\subsection{Критерии классификации}

Выбор критериев классификации обусловлен необходимостью комплексной оценки методов обнаружения признаков генерации. Уровень декомпозиции изображения определяет подход к анализу визуальной информации, вычислительная сложность характеризует практическую применимость метода, а устойчивость к шумам отражает его надёжность при работе с данными, подверженными искажениям. Критерии классификации и классы, соответствующие критериям представлены в таблицах~\ref{tab:classification-criteria}-\ref{tab:noise-resistance-classification}.

\begin{table}[H]
	\centering
	\caption{Критерии классификации методов обнаружения признаков генерации на изображении}
	\label{tab:classification-criteria}
	\begin{tabularx}{\textwidth}{|>{\centering\arraybackslash}p{4cm}|Y|}
		\hline
		\textbf{Критерий} & \textbf{Описание} \\
		\hline
		Уровень декомпозиции изображения & Уровень признаков, на которых фокусируется метод при анализе изображения \\
		\hline
		Вычислительная сложность & Количество параметров метода \\
		\hline
		Устойчивость к шумам & Способность метода сохранять точность при работе с зашумлёнными изображениями \\
		\hline
	\end{tabularx}
\end{table}

\begin{table}[H]
	\centering
	\caption{Классификация по уровню декомпозиции изображения}
	\label{tab:analysis-level-classification}
	\begin{tabularx}{\textwidth}{|>{\centering\arraybackslash}p{4cm}|Y|}
		\hline
		\textbf{Класс} & \textbf{Описание} \\
		\hline
		Микроскопические & Анализ на уровне интенсивностей пикселей \\
		\hline
		Мезоскопические & Анализ признаков промежуточного уровня между пикселями и семантическими объектами \\
		\hline
		Макроскопические & Анализ глобальных признаков, семантических связей и пространственных отношений между частями изображения \\
		\hline
		Холистические & Анализ изображения как целостной системы с учетом четырех взаимосвязей всех элементов и глобального контекста \\
		\hline
	\end{tabularx}
\end{table}

Выбор количественных границ для классификации по вычислительной сложности обусловлен распределением параметров рассмотренных методов. Значение 500 тысяч определено как верхняя граница компактных архитектур, к которым относятся свёрточные сети с объёмом 28-30 тысяч параметров, пригодные для устройств с ограниченными ресурсами. 
\begin{table}[h]
	\centering
	\caption{Классификация по вычислительной сложности}
	\label{tab:complexity-classification}
	\begin{tabularx}{\textwidth}{|>{\centering\arraybackslash}p{4cm}|Y|}
		\hline
		\textbf{Класс} & \textbf{Описание} \\
		\hline
		Компактные & Количество параметров метода меньше 500\,000 \\
		\hline
		Средние & Количество параметров метода в интервале от 500\,000 до 10\,000\,000\\
		\hline
		Крупные &  Количество параметров метода больше 10\,000\,000 \\
		\hline
	\end{tabularx}
\end{table}
Нижняя граница крупных методов установлена на 10 миллионов, поскольку трансформеры и капсульные сети, обеспечивающие глобальный анализ, имеют от 70 до 80 миллионов параметров и выше.

\begin{table}[H]
	\centering
	\caption{Классификация по устойчивости к шумам}
	\label{tab:noise-resistance-classification}
	\begin{tabularx}{\textwidth}{|>{\centering\arraybackslash}p{4cm}|Y|}
		\hline
		\textbf{Класс} & \textbf{Описание} \\
		\hline
		Неустойчивые & Методы, точность которых значительно снижается при работе с зашумлёнными изображениями \\
		\hline
		Устойчивые & Методы, точность которых незначительно снижается при работе с зашумлёнными изображениями \\
		\hline
	\end{tabularx}
\end{table}


\subsection{Результаты классификации}

Результаты классификации обозреваемых решений по вышеописанным критериям представлены в таблице~\ref{tab:methods-classification}.

\begin{table}[H]
	\centering
	\caption{Классификация методов обнаружения признаков генерации на изображении}
	\label{tab:methods-classification}
	\begin{tabularx}{\textwidth}{|>{\centering\arraybackslash}p{3.4cm}|Y|>{\centering\arraybackslash}p{3.4cm}|Y|}
		\hline
		\textbf{Метод} & \textbf{Уровень декомпозиции изображения} & \textbf{Вычисли-тельная сложность} & \textbf{Устойчивость к шумам} \\
		\hline
		Meso-4 & Мезоскопические & Компактные & Неустойчивые \\
		\hline
		Meso Inception-4 & Мезоскопические & Компактные & Устойчивые \\
		\hline
		Capsule Forensics & Макроскопические & Средние & Неустойчивые \\
		\hline
		Vision Transformer & Холистические & Крупные & Устойчивые \\
		\hline
	\end{tabularx}
\end{table}

\section{Сравнительный анализ рассмотренных решений}
\label{sec:comparative-analysis}

В данном разделе представляются критерии и проведение сравнения рассмотренных решений.

\subsection{Критерии сравнения}

Критерии сравнения представлены в таблице~\ref{tab:compare-criteria}.

\begin{table}[H]
	\centering
	\caption{Критерии сравнения методов обнаружения признаков генерации}
	\label{tab:compare-criteria}
	\begin{tabularx}{\textwidth}{|>{\centering\arraybackslash}p{5cm}|Y|}
		\hline
		\textbf{Критерий} & \textbf{Описание} \\
		\hline
		Точность на наборе данных FaceForensics & Доля верно классифицированных изображений на наборе FaceForensics \\
		\hline
		Количество параметров & Количество параметров метода \\
		\hline
		Пороговое значение PSNR & Минимальное значение PSNR, при котором снижение точности метода не превышает 1\% \\
		\hline
	\end{tabularx}
\end{table}

\subsection{Результаты сравнения}

Результаты сравнительного анализа рассмотренных решений представлены в таблице~\ref{tab:methods-compare}.

\begin{table}[H]
	\centering
	\caption{Сравнительный анализ методов обнаружения признаков генерации}
	\label{tab:methods-compare}
	\begin{tabularx}{\textwidth}{|>{\centering\arraybackslash}p{3.4cm}|Y|Y|>{\centering\arraybackslash}p{3.5cm}|}
		\hline
		\textbf{Метод} & 
		\textbf{Точность на FaceForensics, \%} &
		\textbf{Количество параметров, шт.} &
		\textbf{Пороговое значение PSNR, дБ.} \\
		\hline
		Meso-4 & 94.6 & 27\,997 & 30 \\
		\hline
		Meso Inception-4 & 96.8 & 28\,615 & 28 \\
		\hline
		Capsule Forensics & 99.37 & 2\,800\,000 & 34 \\
		\hline
		Vision Transformer & 97.0 & 84\,000\,000 & 24 \\
		\hline
	\end{tabularx}
\end{table}

Из результатов сравнительного анализа можно сделать вывод, что капсульные ИНС демонстрируют наивысшую точность на незашумлённых изображениях, в то время как свёрточные ИНС демонстрируют относительно высокую точность на изображениях любого вида. Трансформеры обладают наибольшей устойчивостью к шуму, что обеспечивается наибольшим количеством параметров.

\section{Формализация изображения}

\textit{Цветом} называется вектор $\vec{c}$, описывающий восприятие электромагнитного излучения в видимом диапазоне и характеризующийся набором интенсивностей, каждая из которых соответствует отдельному каналу восприятия:
\begin{equation}
	\vec{c} = (c_0, c_1, \ldots, c_{N-1}) \in \mathbb{R}^{N}_+,
\end{equation}
где $N \in \mathbb{N}$ -- количество каналов восприятия. Множество всех возможных цветов c $N$ каналами восприятия называется \textit{цветовым пространством} и обозначается $\mathbb{C}_N \subseteq \mathbb{R}^N_+$.

\textit{Изображением} $I$ называется отображение, сопоставляющее каждой точке некоторого координатного пространства $\Omega \subset \mathbb{R}^M_+$ цвет в выбранном цветовом пространстве $\mathbb{C}_N$:
\begin{equation}
	I: \Omega \rightarrow \mathbb{C}_N
\end{equation}
При этом выделяются следующие виды изображений:
\begin{itemize}
	\item изображение называется \textit{непрерывным}, если его область определения $\Omega$ является связным подмножеством $\mathbb{R}^M_+$;
	\item изображение называется \textit{дискретным (цифровым)}, если его область определения $\Omega$ является конечным подмножеством $\mathbb{Z}^M_+$.
\end{itemize}

\textit{Тензором} некоторого дискретного изображения $I$ называется тензор $T^I \in \mathbb{R}^{L_1 \times L_2 \times \cdots \times L_M \times N}_+$, элементы которого определяются следующим правилом:
\begin{equation}
	T^I_{i_1, i_2, \ldots, i_M, k} = I^{(k)}(i_1, i_2, \ldots, i_M), \quad i_\mu = \overline{0, L_\mu-1}, \quad k = \overline{0, N-1}
\end{equation}
где
\begin{itemize}
	\item $I^{(k)}(i_1,\ldots,i_M)$ -- значение интенсивности сигнала $k$-го канала восприятия цвета в точке с координатами $(i_1,\ldots,i_M)$;
	\item $L_\mu \in \mathbb{N}$ -- размер изображения по $\mu$-му координатному направлению;
	\item $N \in \mathbb{N}$ -- количество каналов восприятия.
\end{itemize}

\section{Формализация задачи}
\label{sec:task-formalization}

Пусть $I$ является некоторым цифровым двумерным изображением, а событие $A~=~\{\textit{изображение содержит признаки генерации}\}$. \textit{Функцией дискриминации}~$F_{\tau}$ называется следующая функция:
\begin{equation}
	F_{\tau}(I) = 
	\begin{cases}
		1, & P({A \mid I}) \geqslant \tau \\
		0, & \textit{иначе}
	\end{cases},
\end{equation}
где
\begin{itemize}[label=---]
	\item $\tau \in (0;1)$ -- пороговая вероятность;
	\item $P({A \mid I})$ -- вероятность наличия признаков генерации на изображении $I$.
\end{itemize}
Возвращаемые значения функции дискриминации $F_{\tau}(I)$ отвечают следующему правилу:
\begin{equation}
	\begin{cases}
		1, & \textit{изображение $I$ содержит признаки генерации} \\
		0, & \textit{изображение $I$ не содержит признаки генерации} 
	\end{cases}
\end{equation}
При этом следует отметить, что наличие признаков генерации является необходимым, но не достаточным условием для отнесения изображения к категории <<поддельных>>, так как подделка предполагает также целенаправленный характер фальсификации. Однако в рамках формализации задачи \textbf{допускается} интерпретация факта наличия признаков генерации на изображении как индикатора его поддельности, то есть возвращаемые значения функции дискриминации также понимаются следующим образом:
\begin{equation}
	\begin{cases}
		1, & \textit{изображение $I$ является поддельным} \\
		0, & \textit{изображение $I$ не является поддельным} 
	\end{cases}
\end{equation}

Задача обнаружения признаков генерации на изображении состоит в том, чтобы определить апостериорную вероятность $P(A \mid I)$. Функциональная постановка задачи представлена на рисунке~\ref{img:fpz}.

\includeimage
{fpz}
{f} % Обтекание (без обтекания)
{H}
{\textwidth} % Ширина рисунка
{Функциональная постановка задачи обнаружения признаков генерации на изображении}

\section{Выводы из аналитического раздела}

В результате проведенного анализа для реализации метода обнаружения признаков генерации на изображении с использованием нейронной сети были решены следующие задачи:
\begin{itemize}
	\item проанализирована предметная область;
	\item проведен сравнительный анализ известных нейросетевых решений: сверточных ИНС, капсульных ИНС и трансформеров;
	\item предложены и обоснованы критерии классификации рассмотренных решений.
\end{itemize}	

В результате проведенных классификации и сравнительного анализа было установлено, что:
\begin{itemize}	
	\item вычислительная сложность метода напрямую связана с уровнем декомпозиции изображения;
	\item капсульные ИНС достигают наибольшей точности на наборе FaceForensics, свёрточные ИНС обеспечивают баланс между точностью и устойчивостью к шуму, а трансформеры демонстрируют наибольшую устойчивость к искажениям данных, имея наибольшее количество параметров.
\end{itemize}

Предложена формализация задачи в виде функциональной модели, включающая определение апостериорной вероятности наличия признаков генерации и формальное задание функции дискриминации, что позволяет  интерпретировать задачу обнаружения признаков генерации как задачу бинарной классификации.